\chapter{Priorités opératoires}\label{ch:g7:priorities}

Que vaut $3 + 4 \times 5$~? En calculant de gauche à droite on
obtient $35$~; en multipliant d'abord on obtient $23$. Les
mathématiques n'admettent qu'une seule réponse, d'où des règles de
circulation pour les calculs --- les priorités. Ce chapitre les
énonce et introduit une première identité algébrique, la
distributivité.

\section{Les règles de priorité}

\begin{theorem}[Règles de priorité]\label{thm:g7:priorities:rules}
Dans un calcul sans parenthèses~:
\begin{enumerate}
  \item les multiplications et divisions s'effectuent \emph{avant}
        les additions et soustractions~;
  \item les opérations de même niveau s'effectuent de gauche à
        droite.
\end{enumerate}
Les parenthèses l'emportent sur tout~: ce qui est entre parenthèses
se calcule d'abord, en commençant par les plus intérieures.
\end{theorem}

\begin{example}\label{ex:g7:priorities:basic}
Une opération par ligne, en soulignant (mentalement) ce que l'on
calcule~:
\begin{align*}
3 + 4 \times 5 &= 3 + 20 = 23
&&\text{(multiplication d'abord)}\\
20 - 8 + 2 &= 12 + 2 = 14
&&\text{(même niveau~: gauche à droite)}\\
18 \div 3 \times 2 &= 6 \times 2 = 12
&&\text{(même niveau~: gauche à droite)}\\
(3 + 4) \times 5 &= 7 \times 5 = 35
&&\text{(parenthèses d'abord).}
\end{align*}
Attention aux pièges~: $20 - 8 + 2$ n'est \emph{pas} $20 - 10$, et
$18 \div 3 \times 2$ n'est \emph{pas} $18 \div 6$.
\end{example}

\begin{example}[Parenthèses imbriquées]\label{ex:g7:priorities:nested}
\begin{align*}
50 - 2 \times \bigl(3 + (10 - 4)\bigr)
&= 50 - 2 \times (3 + 6) && \text{(parenthèse intérieure)}\\
&= 50 - 2 \times 9 && \text{(parenthèse restante)}\\
&= 50 - 18 && \text{(multiplication)}\\
&= 32 .
\end{align*}
\end{example}

\begin{remark}[Les barres de fraction cachent des parenthèses]
Une barre de \omterm{def:g6:fractions:def}{fraction} regroupe son \omterm{def:g4:fractions:def}{numérateur} et son \omterm{def:g4:fractions:def}{dénominateur},
comme s'il y avait des parenthèses invisibles~:
\[
\frac{12 + 8}{4} = \frac{20}{4} = 5,
\qquad\text{et non}\quad 12 + \frac{8}{4} = 14 .
\]
\end{remark}

\begin{method}[Calculer en sécurité]\label{met:g7:priorities:safe}
\begin{enumerate}
  \item Repérer l'opération à faire en premier (parenthèses les plus
        intérieures, puis $\times$ et $\div$, puis $+$ et $-$)~;
  \item calculer \emph{uniquement celle-là}, et réécrire toute
        l'expression avec le résultat à la place~;
  \item recommencer jusqu'à ce qu'il reste un seul nombre --- une
        opération par ligne, sans raccourci.
\end{enumerate}
\end{method}

\section{Distributivité}

\begin{theorem}[Distributivité]\label{thm:g7:priorities:distributivity}
Pour tous nombres $k$, $a$, $b$~:
\[
k \times (a + b) = k \times a + k \times b,
\qquad
k \times (a - b) = k \times a - k \times b .
\]
\end{theorem}

\begin{proof}[Preuve par une figure]
Compter l'\omterm{def:g6:measure:area}{aire} d'un rectangle de largeur $k$ découpé en deux
rectangles plus petits de longueurs $a$ et $b$~: en un seul morceau
c'est $k \times (a+b)$~; en deux parties c'est
$k \times a + k \times b$. Même rectangle, même \omterm{def:g6:measure:area}{aire}.
\end{proof}

\begin{omfigure}
\begin{tikzpicture}[scale=0.75]
  \fill[omDef!14] (0,0) rectangle (4.2,1.8);
  \fill[omThm!14] (4.2,0) rectangle (6.4,1.8);
  \draw[thick] (0,0) rectangle (6.4,1.8);
  \draw[thick] (4.2,0) -- (4.2,1.8);
  \node[below, font=\small] at (2.1,0) {$a$};
  \node[below, font=\small] at (5.3,0) {$b$};
  \node[left, font=\small] at (0,0.9) {$k$};
  \node[font=\small] at (2.1,0.9) {$k \times a$};
  \node[font=\small] at (5.3,0.9) {$k \times b$};
\end{tikzpicture}

{\small Distributivité vue comme \omterm{def:g6:measure:area}{aires}~: le grand rectangle
$k \times (a + b)$ est la \omterm{def:g1:addition:def}{somme} des deux plus petits.}
\end{omfigure}

\begin{example}[Les deux sens]\label{ex:g7:priorities:distributivity}
\emph{Développer} (de gauche à droite) facilite le calcul mental~:
\[
7 \times 103 = 7 \times (100 + 3) = 700 + 21 = 721,
\qquad
6 \times 98 = 6 \times (100 - 2) = 600 - 12 = 588 .
\]
\emph{Factoriser} (de droite à gauche) regroupe un facteur commun~:
\[
23 \times 12 + 23 \times 8 = 23 \times (12 + 8) = 23 \times 20 = 460 .
\]
\end{example}

\begin{example}[Traduire des mots en calcul]\label{ex:g7:priorities:words}
«~La \omterm{def:g1:addition:def}{somme} de $8$ et du \omterm{def:g2:mult:def}{produit} de $5$ par $3$~» s'écrit
$8 + 5 \times 3 = 23$. «~Le \omterm{def:g2:mult:def}{produit} de $5$ par la \omterm{def:g1:addition:def}{somme} de $8$ et
$3$~» s'écrit $5 \times (8 + 3) = 55$~: les mots «~le \omterm{def:g2:mult:def}{produit}
de~\dots\ par la \omterm{def:g1:addition:def}{somme} de~\dots~» placent des parenthèses autour de
la \omterm{def:g1:addition:def}{somme}.
\end{example}

\section{Exercices}

\begin{exercise}[$\star$]\label{exo:g7:priorities:1}
Calculer, une opération par ligne~:
\[
5 + 3 \times 6, \qquad
(5 + 3) \times 6, \qquad
30 - 12 \div 4, \qquad
(30 - 12) \div 4 .
\]
\end{exercise}

\begin{exercise}[$\star$]\label{exo:g7:priorities:2}
Calculer~:
\[
24 - 6 + 2, \qquad
24 \div 6 \times 2, \qquad
24 - 6 \times 2, \qquad
24 \div (6 \times 2).
\]
\end{exercise}

\begin{exercise}[$\star$]\label{exo:g7:priorities:3}
Calculer étape par étape~:
\[
7 + 3 \times (10 - 6), \qquad
40 - (2 + 3) \times 4, \qquad
60 \div \bigl(2 \times (7 - 4)\bigr).
\]
\end{exercise}

\begin{exercise}[$\star$]\label{exo:g7:priorities:4}
Calculer les \omterm{def:g6:fractions:def}{fractions} (parenthèses invisibles~!)~:
\[
\frac{18 + 6}{8}, \qquad
\frac{30}{7 - 2}, \qquad
\frac{5 \times 8}{4 + 6} .
\]
\end{exercise}

\begin{exercise}[$\star$]\label{exo:g7:priorities:5}
Utiliser la distributivité pour calculer mentalement~:
\[
8 \times 102, \qquad
25 \times 99, \qquad
14 \times 12 + 14 \times 8, \qquad
37 \times 45 - 37 \times 35 .
\]
\end{exercise}

\begin{exercise}[$\star$]\label{exo:g7:priorities:6}
Écrire le calcul en une ligne, avec des parenthèses si besoin, puis
calculer~: «~la \omterm{ex:g1:subtraction:difference}{différence} de $50$ et du \omterm{def:g2:mult:def}{produit} de $6$ par $7$~»~;
«~le \omterm{def:g2:mult:def}{produit} de la \omterm{def:g1:addition:def}{somme} de $4$ et $5$ par la \omterm{ex:g1:subtraction:difference}{différence} de $10$ et
$8$~».
\end{exercise}

\begin{exercise}[$\star$]\label{exo:g7:priorities:7}
Amir a acheté $3$ cahiers à $2.50$ chacun et $2$ stylos à $1.20$
chacun. Écrire son total en une seule expression, puis le calculer.
\end{exercise}

\begin{exercise}[$\star\star$]\label{exo:g7:priorities:8}
Placer des parenthèses dans $4 + 2 \times 5 - 1$ pour obtenir~:
(a)~$13$~; (b)~$29$~; (c)~$24$. (L'un des trois n'a besoin d'aucune
parenthèse --- lequel~?)
\end{exercise}

\begin{exercise}[$\star\star$]\label{exo:g7:priorities:9}
Les expressions $12 + 4 \times k$ et $16 \times k$ sont-elles les
mêmes pour tout nombre $k$~? Tester avec $k = 1$ et $k = 2$, et
conclure. Quelle erreur de priorité ferait passer de l'une à
l'autre~?
\end{exercise}

\begin{exercise}[$\star\star$]\label{exo:g7:priorities:10}
Un billet de cinéma coûte $9$ euros, mais un groupe de $n$ personnes
paie un forfait de $12$ euros plus $6$ euros par personne.
\begin{enumerate}
  \item Écrire une expression pour le prix du groupe, et la calculer
        pour $n = 5$.
  \item À partir de combien de personnes le tarif groupe est-il plus
        avantageux que $n$ billets individuels~? (Essayer quelques
        valeurs de $n$.)
\end{enumerate}
\end{exercise}

\begin{exercise}[$\star\star\star$]\label{exo:g7:priorities:11}
En utilisant chacun des chiffres $1$, $2$, $3$, $4$ exactement une
fois, les quatre opérations et des parenthèses à volonté, écrire des
expressions égales à $24$, à $10$, et à $1$.
\end{exercise}
